% Source: https://artofproblemsolving.com/community/c3h570919_2013_national_team_10
% Stable identity: aops-topic-570919-q10
With $m=n+k$: $\dfrac{n(n+1)}{2}=\dfrac{m(m+1)}{2}-\dfrac{n(n+1)}{2}$, so $m(m+1)=2n(n+1)$. Multiplying by $4$ and adding $1$: $(2m+1)^2-2(2n+1)^2=-1$, a Pell equation $X^2-2Y^2=-1$ with $X=2m+1$, $Y=2n+1$. Its solutions are $(X,Y)=(1,1),(7,5),(41,29),\dots$, giving $(n,m)=(2,3),(14,20),(84,119),\dots$. (Check: $1+2=3$, and $1+2+\dots+14=105=15+16+\dots+20$.) With $n<20$ the largest is $n=14$, $k=20-14=\boxed{6}$.
